\documentclass[11pt]{article}

\usepackage[margin=1in]{geometry}
\usepackage{amsmath,amssymb}
\usepackage{enumitem}

\setlength{\parindent}{0pt}
\setlength{\parskip}{0.6em}

\begin{document}

\begin{center}
    {\Large \textbf{MATH 4790: Quiz 4}}\\[0.3em]
    March 30, 2026
\end{center}

\vspace{1em}

\textbf{Name:} \hfill \textbf{Score:} \rule{1.5cm}{0.4pt}/15

\vspace{1em}

\hrule
\vspace{1em}

\textbf{Instructions:} Show all work and justify your answers. Only calculators approved for Exam~FM are permitted. Credit will not be given to just guessing the correct answer.

\vspace{1em}

% -------------------- QUESTION 1 (OLD Q3) --------------------

\noindent\textbf{Question 1.}  

A loan is repaid with 18 annual level payments made at the end of each year. The annual effective interest rate is $5.4\%$.

The interest portion of the second payment is equal to the principal portion of the $t$th payment.

Find $t$.

\medskip

\begin{center}

\textbf{(A)} $7$
\hspace{1cm}
\textbf{(B)} $8$
\hspace{1cm}
\textbf{(C)} $9$
\hspace{1cm}
\textbf{(D)} $10$
\hspace{1cm}
\textbf{(E)} $11$

\end{center}

\newpage

% -------------------- QUESTION 2 --------------------

\noindent\textbf{Question 2.}  

An office takes out a \$1{,}250 equipment loan at time \(t=0\), where time is measured in years. The annual effective interest rate is \(20\%\) for the first two years, and \(25\%\) for all years thereafter.

The loan is scheduled to be repaid using the following payments:
\begin{itemize}[leftmargin=2em]
    \item a payment of \(X\) at time \(t=2\),
    \item a payment of \(Y\) at time \(t=4\), and
    \item a final payment of \(625\) at time \(t=5\).
\end{itemize}

No other payments are made.

At time \(t=3\), the outstanding balance on the loan is \(1{,}000\).

Find \(\dfrac{X}{Y}\).

\medskip

\begin{center}

\textbf{(A)} \(\dfrac{5}{4}\)
\hspace{1cm}
\textbf{(B)} \(\dfrac{4}{3}\)
\hspace{1cm}
\textbf{(C)} \(\dfrac{3}{2}\)
\hspace{1cm}
\textbf{(D)} \(2\)
\hspace{1cm}
\textbf{(E)} \(\dfrac{9}{4}\)

\end{center}

\newpage

% -------------------- QUESTION 3 (OLD Q1) --------------------

\noindent\textbf{Question 4.}  

A $20$-year loan is repaid with end-of-month payments. Let \(j\) denote the nominal annual interest rate compounded monthly.

For the first year, payments are \$300 each. The remaining payments are \$500 each.

\textit{\underline{Right before}} the 12th payment, the outstanding loan balance is \$28{,}000.

Find \(j\).

\medskip

\begin{center}

\textbf{(A)} \(19.5\%\)
\hspace{1cm}
\textbf{(B)} \(20.1\%\)
\hspace{1cm}
\textbf{(C)} \(20.7\%\)
\hspace{1cm}
\textbf{(D)} \(21.3\%\)
\hspace{1cm}
\textbf{(E)} \(21.9\%\)

\end{center}

\end{document}
